\section{把真实机器人数据提升为 4D 并预训练}
\label{app:pretraining}

\subsection{数据、划分与时间}
\label{app:droid}
预训练语料是 DROID \cite{droid}，以 15\,Hz 录制，含两台外部相机与一台腕部相机，并配有 PointWorld 为其发布的 3D 点轨迹与优化后的相机标定 \cite{pointworld}。我们要求关节序列、末端位姿、夹爪状态、指令与相机身份完全一致，才把每个发布的片段与其原始 DROID 回合关联起来；采集时钟不单调或关联有歧义的片段被排除。模型看到的是第一台外部相机。训练单元是一个因果观测窗口及其对齐的未来与动作标签。我们在任何筛选之前，按原始标识把回合分到训练、开发与留出划分（90/5/5，固定随机种子），因此同一回合的重叠片段不会跨越划分。表~\ref{tab:pretraining_data} 报告每个阶段的产出，并分别统计回合、窗口、动作与有效标签。

% Planned results: DROID lifting statistics and stage yield (App. E). PW0 manifest facts for reference (confirm before submission): 42,935 released episodes; 42,087 trainable; 598,889 main-camera windows (train 37,874 ep. / 538,740 win.; val 4,213 / 60,149); window action validity 96.88%.
\begin{table}[t]
\centering\footnotesize
\caption{\tablelang{\textbf{DROID lifted to 4D.} Yield of each stage. Windows are causal observation windows with aligned
labels. Exogenous windows show target motion before the first robot--target contact (Appendix~\ref{app:lifting}). The last
two rows compare all windows with valid windows.}{\textbf{提升为 4D 的 DROID。}各阶段的保留量。窗口指带有对齐标签的因果观测窗口。
外源窗口指在机器人与目标首次接触前目标已在运动的窗口（附录~\ref{app:lifting}）；最后两行比较全部窗口与有效窗口。}}
\label{tab:pretraining_data}
\setlength{\tabcolsep}{3.4pt}
\begin{tabular}{@{}lccc@{}}
\toprule
 & \tablelang{Train}{训练} & \tablelang{Dev.}{开发} & \tablelang{Held-out}{留出} \\
\midrule
\tablelang{Episodes with PointWorld tracks}{含 PointWorld 轨迹的 episode} & \resultplaceholder & \resultplaceholder & \resultplaceholder \\
\quad\tablelang{associated to raw DROID}{与原始 DROID 对齐} & \resultplaceholder & \resultplaceholder & \resultplaceholder \\
\quad\tablelang{one rigid target named and boxed}{指令只指定一个刚体且已检测} & \resultplaceholder & \resultplaceholder & \resultplaceholder \\
\tablelang{Windows}{窗口} & \resultplaceholder & \resultplaceholder & \resultplaceholder \\
\quad\tablelang{exogenous (\%)}{外源（\%）} & \resultplaceholder & \resultplaceholder & \resultplaceholder \\
\quad\tablelang{with causal pose input}{含因果位姿输入} & \resultplaceholder & \resultplaceholder & \resultplaceholder \\
\tablelang{Valid actions (\%)}{有效动作（\%）} & \resultplaceholder & \resultplaceholder & \resultplaceholder \\
\tablelang{Valid SE(3) labels (horizons, \%)}{有效 SE(3) 标签（时域，\%）} & \resultplaceholder & \resultplaceholder & \resultplaceholder \\
\tablelang{Valid 3D flow labels (points, \%)}{有效 3D 流标签（点，\%）} & \resultplaceholder & \resultplaceholder & \resultplaceholder \\
\tablelang{Speed, all / valid (cm/s)}{速度，全部 / 有效（cm/s）} & \resultplaceholder[0.0]/\resultplaceholder[0.0] & \resultplaceholder[0.0]/\resultplaceholder[0.0] & \resultplaceholder[0.0]/\resultplaceholder[0.0] \\
\tablelang{Occluded, all / valid (\%)}{遮挡，全部 / 有效（\%）} & \resultplaceholder[00]/\resultplaceholder[00] & \resultplaceholder[00]/\resultplaceholder[00] & \resultplaceholder[00]/\resultplaceholder[00] \\
\bottomrule
\end{tabular}
\end{table}


\paragraph{时间。}所有时间都是采集时间戳，我们不把帧重新标注为名义帧率。历史偏移位于 15\,Hz 帧不会恰好落在的 50\,ms 网格上，因此我们与部署时完全一样（附录~\ref{app:tokens}），在查询时刻之前已输出的估计之间对因果输入在这些偏移处重采样，不做外推。目标标签以同样的方式对离线轨迹插值（位置线性插值，旋转沿测地线插值），不跨越超过两个源周期的空缺，因此它们位于主模型的时域上。流标签保留 PointWorld 自己的记录时刻（前方 \fillin{0.13} 到 \fillin{1.13}\,s），不做插值。世界读出头嵌入每个时域在物理时间中的偏移 $\delta_h$（附录~\ref{app:world}），因此场景查询在预训练时在这些偏移上受到监督，之后在 50\,ms 网格上受到监督。

\paragraph{动作。}我们在经过标定的法兰到工具中心点变换之后，在机器人基座坐标系中表示机器人状态与笛卡尔指令。指令标签遵循附录~\ref{app:action}：每个标签是从为前一个生效时刻指令的位姿到为当前生效时刻指令的位姿的增量，DROID 的指令位姿被插值到所采样延迟对应的生效时刻 $\tau_t+\ell+j\Delta_a$。数值有限且其参考状态可因果获得的标签才被接受。机器人的执行限幅不是接受规则，因为每条指令每轴 5\,cm 的限幅会丢弃运动最快四分位中 \fillin{XX}\% 的指令，使基座偏向慢速示范。我们只屏蔽 DROID 采集器没有发出指令的步（窗口槽位的 \fillin{3.1}\%）。动作与流的有效性相互独立，因此无效的动作不会去掉其窗口的流标签。

\subsection{目标及其位姿}
\label{app:droid_target}
DROID 不提供物体位姿，其指令只说出被操作物体的名字而不标出它。我们把指令中提到的物体作为目标，即依存句法分析器提取的主要动词的直接宾语（``put the mug in the sink'' 中的 ``the mug''），每个回合一个持续存在的实例。Grounding DINO \cite{groundingdino} 在该短语得分首次达到 0.35 的帧中给出包围盒，SAM~2.1 \cite{sam2} 跟踪其掩码。对每个掩码非空的帧，我们以紧致 2D 包围盒、RGB 帧与固定内参提示 WildDet3D \cite{wilddet3d}。它从单张图像、用自身的几何估计而不借助外部深度预测度量 3D 包围盒，PointWorld 优化后的外参把包围盒映射到基座坐标系。当两台外部相机给出的包围盒中心在基座坐标系中相距 5\,cm 以内时，我们接受该包围盒。指令中没有物体、有多个物体或物体非刚性的回合，以及没有被接受的包围盒的回合，都被排除。空掩码记为缺失观测，我们不用上一个位姿填补。

\paragraph{物体坐标系。}3D 包围盒只在其对称性范围内确定坐标系，检测器在相邻帧之间可能排列或翻转其坐标轴。因此我们把每个新包围盒映射到与上一个有效估计最接近的等价坐标系（保持其尺寸的坐标轴排列与符号翻转），从而为每个回合固定一个物体坐标系。经过这一步后仍然存在的相邻帧间超过 $45^\circ$ 的旋转跳变（占 \fillin{X.X}\% 的帧）被标为无效，而不是被平滑。与仿真中一样，对称目标绕对称轴的旋转被排除在损失之外（第~\mainref{sec:predict}~节）。

\subsection{因果输入与离线标签}
\label{app:lifting}
\begin{figure*}[t]
\centering
\figplaceholder[1.6in]{\textwidth}{图 S1——把真实机器人回合提升为度量 4D（示意）}{%
\textbf{左（共同设置）：}指令及其被加下划线的物体短语（``put \underline{the mug} in the sink''），以及两个外部视角第一帧上的 Grounding DINO 包围盒，其 WildDet3D 包围盒中心在基座坐标系中一致；小插图给出 PointWorld 优化后的相机位姿与原始位姿的对比。\textbf{上行（因果输入路径，\textcolor{obscol}{蓝}）：}帧从左到右到达；SAM~2 向前传播掩码，WildDet3D 把每个带掩码的帧提升为度量 3D 包围盒；输出的位姿画为带可用时刻标记的坐标轴；空掩码显示一个明确的``缺失'' token。没有箭头指向过去。
\textbf{下行（离线标签路径，\textcolor{predcol}{橙}）：}向前与向后传播的掩码，映射到同一物体坐标系的逐帧包围盒，由鲁棒平滑器融合并保留突变事件（抓取）的两条位姿轨迹，50\,ms 时域上的累计 SE(3) 标签 $(\vect d_h,\boldsymbol\phi_h)$；另一分支在每个网格图像块中选取离块中心最近的 PointWorld 轨迹，画出其在 PointWorld 记录时刻的 3D 位移，无效点以灰色显示。竖直虚线标出机器人与目标的首次接触，掩码标签对照在其后去掉世界标签。\textbf{右：}一个训练样本 = 因果输入（上）+ 带掩码的标签（下）。使用一段真实 DROID 帧序列，只用矢量叠加。}
\caption{把 DROID 提升为度量 4D 监督。指令中提到的物体被框出并跟踪。输入以只向前的跟踪重现部署条件，标签则使用整个回合、双向跟踪与平滑，以及 PointWorld 的 3D 轨迹。标签不改写输入，无效标签屏蔽其损失。}
\label{fig:lifting}
\end{figure*}

\paragraph{因果输入。}SAM~2.1（Hiera-small）因果地传播掩码，WildDet3D 以全精度提升每个带掩码的帧，使估计不依赖于 batch 的组成。每次查询只处理新到达的帧，并追加其估计而不改写之前的估计。空掩码、被拒绝的包围盒与旋转跳变把输入有效性置为假；训练使用这些自然出现的缺失状态，不做合成的随机丢弃。对没有 CAD 模型或参考视图的物体，真机使用同一路径（附录~\ref{app:tracker}）。

\paragraph{离线标签。}标签使用整个回合与同一目标短语。SAM~2 向前和向后传播掩码，我们接受两个掩码 IoU 至少为 0.5 的帧。两个方向的包围盒进入同一个鲁棒平滑器：Huber 代价，位姿残差以 5\,mm 与 $3^\circ$ 缩放，速度变化残差权重为 0.1。在空缺处以及超过 0.2\,m/s 或 1\,rad/s 的速度跳变处省略速度项，以保留突变事件。二十次流形 Gauss--Newton 迭代优化该代价，两个包围盒相差超过 20\,mm 或 $15^\circ$ 的帧为无效。目标标签是离线轨迹内的 $\tilde{\vect p}_{t+h}-\tilde{\vect p}_t$ 与 $\Log(\tilde R_{t+h}\tilde R_t^\top)^\vee$。由于它们不是离线未来减去因果估计，因此不会吸收对当前位姿的修正。只有当参考位姿与未来位姿都有效、且身份与物体坐标系相同时，一个时域才有效。

\paragraph{场景流。}PointWorld 的轨迹已经在基座坐标系中表示。对网格中的每个点，标签取其初始投影在该点所在图像块内离该点最近的已发布轨迹，前提是该轨迹初始可见、在图像内有有效的正深度，并与其最近像素处的深度相差在 20\,mm 以内。没有这样的轨迹的图像块为无效，我们不用邻近块填补。标签是轨迹在每个记录时刻相对其初始位置的位移，可见性取自 PointWorld。被遮挡、出图像或深度无效的未来点屏蔽其位移损失，只监督其可见性。所选轨迹到图像块中心的距离中位数为 \fillin{0.5}\,px（第 95 百分位为 \fillin{9}\,px）。

\paragraph{接触与外部运动。}机器人与目标的首次接触是目标包围盒首次进入由机器人运动学得到的夹爪手指几何 1\,cm 范围内的那一帧。若目标在该接触之前、在其 400\,ms 时域内移动超过 1\,cm 或 $5^\circ$，则该窗口为\emph{外部运动}窗口，表~\ref{tab:pretraining_data} 报告其比例。表~\ref{tab:supp_pretraining} 中的掩码标签对照在首次接触之后去掉所有世界标签，并保留动作损失。它保留了预训练低数据增益的 \fillin{XX}\%，因此\fillin{这一增益的大部分并非来自预测机器人如何移动目标}。

\paragraph{有效性偏差。}PointWorld 的轨迹对静态表面的覆盖比对快速运动表面更完整，被接受的包围盒偏向未被遮挡、纹理丰富的目标。表~\ref{tab:pretraining_data} 比较全部窗口与有效窗口之间的目标速度与遮挡情况。表~\ref{tab:supp_pretraining} 中的标签来源对照用从 2D 对应关系 \cite{dpflow} 结合单目度量深度 \cite{moge2} 提升得到的流替代这些轨迹，后者覆盖图像的更大部分。

\subsection{标签质量}
\label{app:label_quality}
% Planned results: controlled DROID-base comparisons and data efficiency (App. E). Simulation post-training from each initialization; DROID held-out forecast error of each base. Real-robot effects of the base: Tables 5, S-real.
\begin{table*}[t]
\centering\small
\caption{\tablelang{\textbf{Controlled comparison of initializations.} Each row post-trains the simulation model from its
own initialization on 10/25/50/100\% of the \bench demonstrations (held-out SR, \%, three post-training seeds). The
generic initialization is the one \ours uses. The last columns give each base's forecast error on held-out DROID episodes
at the last horizon, for target translation (mm) and rotation ($^\circ$) against the smoothed labels and for the flow
endpoint (mm) against PointWorld tracks. Control bases are pretrained at pilot scale against a world--action base at the same
budget (Appendix~\ref{app:scale}).}{\textbf{初始化的受控比较。}每行从各自的初始化出发，用 10/25/50/100\% 的 \bench 示范后训练仿真模型
（留出 SR，\%，三个后训练种子）；通用初始化即 \ours 所用的初始化。最后几列为各基座在留出 DROID episode 上最后一个时域的预测误差：
目标平移（mm）与旋转（$^\circ$，相对平滑标签），以及相对 PointWorld 轨迹的流端点误差（mm）。对照基座按试点规模预训练，并与相同预算的世界—动作基座比较（附录~\ref{app:scale}）。}}
\label{tab:supp_pretraining}
\setlength{\tabcolsep}{4.6pt}
\begin{tabular}{@{}lccccccc@{}}
\toprule
 & \multicolumn{4}{c}{\tablelang{\bench SR, fraction of demonstrations}{\bench SR，示范比例}} & \multicolumn{3}{c}{\tablelang{DROID held-out}{DROID 留出}} \\
\cmidrule(lr){2-5}\cmidrule(lr){6-8}
\tablelang{Initialization}{初始化} & 10\% & 25\% & 50\% & 100\% & mm & $^\circ$ & \tablelang{flow (mm)}{流（mm）} \\
\midrule
\tablelang{Generic visual only (\ours)}{仅通用视觉初始化（\ours）} & \resultplaceholder & \resultplaceholder & \resultplaceholder & \resultplaceholder & -- & -- & -- \\
\tablelang{Action-only base, same episodes}{仅动作基座，相同 episode} & \resultplaceholder & \resultplaceholder & \resultplaceholder & \resultplaceholder & -- & -- & -- \\
\tablelang{World--action base, flow from monocular depth}{世界—动作基座，单目深度提升的流} & \resultplaceholder & \resultplaceholder & \resultplaceholder & \resultplaceholder & \resultplaceholder & \resultplaceholder & \resultplaceholder \\
\tablelang{World--action base, world labels masked after contact}{世界—动作基座，接触后屏蔽世界标签} & \resultplaceholder & \resultplaceholder & \resultplaceholder & \resultplaceholder & \resultplaceholder & \resultplaceholder & \resultplaceholder \\
\tablelang{World--action base, pilot scale}{世界—动作基座，试点规模} & \resultplaceholder & \resultplaceholder & \resultplaceholder & \resultplaceholder & \resultplaceholder & \resultplaceholder & \resultplaceholder \\
\tablelang{World--action base (real-robot base)}{世界—动作基座（真机基座）} & \resultplaceholder & \resultplaceholder & \resultplaceholder & \resultplaceholder & \resultplaceholder & \resultplaceholder & \resultplaceholder \\
\bottomrule
\end{tabular}
\end{table*}

我们在四十个片段上检查标签，其中十四个用于开发、二十六个留出，每个片段五个锚点。这些片段覆盖七个层次：静止目标、平滑平移、旋转、接触或打滑、遮挡与恢复、外观较弱，以及对称或铰接部件。留出片段未用于调整估计器。在对 \fillin{200} 个随机回合的人工检查中，被接受的包围盒在 \fillin{XX}\% 的回合中落在预期目标上。前缀截断与未来替换实验证实，已输出的输入、有效性、身份与参考几何不随未来帧改变。独立的人工标注关键点与多视角几何在约束充分处提供参考，其他地方我们把物理精度标为未验证。在留出片段上，因果路径的平移与旋转误差中位数为 \fillin{X}\,mm 与 \fillin{X}$^\circ$，离线路径为 \fillin{X}\,mm 与 \fillin{X}$^\circ$。身份错误出现在 \fillin{X}\% 的锚点上，被接受窗口中有 \fillin{X}\% 为误接受，标签覆盖 \fillin{XX}\% 的窗口。这次检查用了 \fillin{X} 小时的标注，提升整个 DROID 用了 \fillin{X} GPU 小时。在图像中静止的表面上，所选 PointWorld 轨迹在各时域上的位移为 \fillin{0.9} 与 \fillin{5.8}\,mm（中位数与第 95 百分位），而用单目深度提升的流为 \fillin{18} 与 \fillin{72}\,mm，其尺度随时间漂移。在运动子集上，两种来源都没有参考，我们报告两者与一个前向—后向一致的 2D 跟踪器之间的一致性。

\subsection{训练方案与规模}
\label{app:scale}
\paragraph{仿真模型。}\ours 从通用视觉初始化出发，以表~\ref{tab:hyper} 的方案在 \bench 上训练。我们在相同训练样本数下扫描两种主干规模（参考规模以及固定宽度、加深到约两倍参数的规模）、两种全局 batch（$B_0$ 与其四倍）以及每种 batch 两种学习率（$B_0$ 时为 $\{0.5,1\}\eta_0$，较大 batch 时为 $\{1,2\}\eta_0$），再在每种规模的最佳配置周围做 \fillin{四}次局部细化，全部在开发集场景上选择（图~\ref{fig:scaling}a），以确定其规模与方案。

\paragraph{DROID 基座。}基座使用相同的结构与初始化，只使用真实数据。我们训练了两个试点规模，约为仿真主干的两倍与四倍，在按层抽取的四分之一训练窗口上训练一遍。我们依据留出集上的预测与动作误差，在一块 \cfgGPU{} 的部署预算（p95 出动作时间不超过 200\,ms，且没有超过 250\,ms 的调用）下选择规模。选定的规模（主干参数 \fillin{XX}M）在全部训练窗口上预训练一次，共 \fillin{四}遍，全局 batch 为 \fillin{256}，使用 AdamW、2.5\% 预热以及余弦衰减到峰值学习率的 0.1。开发集损失最低的检查点即为基座（图~\ref{fig:scaling}b）。我们以较小的试点规模与预算训练只用动作、标签来源与掩码标签三个对照，并与同规模的世界—动作基座比较（表~\ref{tab:supp_pretraining}）。

\begin{figure}[t]
\centering
\figplaceholder[1.55in]{\linewidth}{图 S2——训练方案与规模}{%
\textbf{(a) 仿真扫描。}y：开发集 SR；x 分组：两种主干规模；每组内在相同样本量下的四根柱：$(B_0,0.5\eta_0)$、$(B_0,\eta_0)$、$(4B_0,\eta_0)$、$(4B_0,2\eta_0)$；细化结果为小菱形；选定配置以星号标出。
\textbf{(b) DROID 基座。}x：训练样本数（对数）；y：留出集最后一个时域的目标平移误差（左轴，mm）与流端点误差（右轴，mm）；细线为两个试点规模，粗线为主训练；标记选定的检查点；水平点线：同一历史的匀速外推。}
\caption{仿真模型（a）与 DROID 基座（b）的训练方案与规模。}
\label{fig:scaling}
\end{figure}

\subsection{独立的后训练与对照}
\label{app:posttraining}
真机策略从选定的基座初始化，只在传送带示范上微调。基座是在检测器位姿上预训练的，而传送带策略接收基于模型的跟踪器给出的位姿；微调使历史编码器适应跟踪器，我们没有单独测量这一位姿来源变化的代价。仿真权重不用于初始化任何真机策略，基座也不接触仿真数据。每个策略创建自己的优化器、学习率计划、仅用于训练的归一化与读出头，不兼容的动作或状态投影在每个对照中都以相同种子重新初始化。不使用基座的参考策略从通用视觉初始化出发，在相同示范上微调，在开发试验上进行相同的搜索（\fillin{三}种学习率与两种训练长度），并采用相同的评测。仿真中的受控比较以下游方案、与通用初始化相同的数据比例在 \bench 上对基座做后训练。只用动作的对照在相同窗口上预训练 R0，因此其世界查询没有目标。标签来源对照只替换流标签，掩码标签对照在首次接触后去掉世界标签（表~\ref{tab:supp_pretraining}）。
